\documentclass[twocolumn]{iopjournal}

\usepackage{amsmath,amssymb,amsfonts}

\fancyhfoffset[L]{0mm}
\setlength{\marginparwidth}{0pt}
\setlength{\marginparsep}{0pt}

% ==========================================
% ZENODO PREPRINT HEADERS (ACTIVE)
% ==========================================
\renewcommand{\articletype}[1]{%
  {\vspace*{-8mm}\noindent \Large \sf Preprint}\par
  \vspace*{6mm}\noindent{\scriptsize \sf{\bfseries \MakeUppercase{#1}}}}

\fancyhead[L]{}
\fancyhead[R]{}

\begin{document}

\articletype{Letter}

\title{The Gravitational Constitutive Law\\
from the Tetrad Frame Geometry}

\author{Stephen Nelson}

\affil{Field Dynamics Research}

\ead{stephen@fielddynamics.org}

\begin{abstract}
The AQUAL field equation requires an interpolating function
$\mu(x)$ whose form has remained undetermined since Milgrom's
original formulation in 1983. We derive it through three
methodologically distinct routes. The first reads the graded
rank structure of the metric and tetrad equation. The second
derives the function from equilibrium field closure within the
AQUAL framework. The third, resting on the companion
letter~[Nelson 2026a], expresses the exchange symmetry of the
tetrad equation's two parity-related bounded regions as a
functional equation on the odds ratio
$\Omega = \mu/(1-\mu)$, forcing $\Omega(x) = x$ and therefore
$\mu(x) = x/(x+1)$. All three routes, through different
mathematics, identify the same algebraic fact: the capacity
ratio between the two tetrahedral sectors equals the
dimensionless acceleration. The constitutive law is the
odds-to-fraction normalisation of that ratio. Its consequences
include the complement symmetry $\mu(x) + \mu(1/x) = 1$ and
the baryonic Tully-Fisher relation as a theorem.
\end{abstract}

% ==================================================================

\section{The open problem}

In 1983, Milgrom~\cite{Milgrom1983} showed that the dynamics of
galaxies at low acceleration can be captured by a single
modification to the Poisson equation, provided an interpolating
function $\mu(x)$ mediates between the Newtonian regime at high
acceleration and a scale-invariant regime at low acceleration.
Bekenstein and Milgrom~\cite{BekensteinMilgrom1984} subsequently
embedded this modification in the AQUAL field equation,
\begin{equation}
  \nabla \cdot \!\left[\mu\!\left(\frac{|\nabla\Phi|}{a_0}\right)
  \nabla\Phi\right] = 4\pi G\rho,
  \label{eq:aqual}
\end{equation}
which derives from an action principle and therefore conserves
energy, momentum, and angular momentum. The function $\mu(x)$,
where $x = |\nabla\Phi|/a_0$, must approach unity at large $x$
to recover Newtonian gravity and must scale linearly with $x$
as $x \to 0$ to produce flat rotation curves.

Infinitely many functions satisfy both limits. Among the
candidates explored in the literature, the ``simple'' form
$x/(x+1)$~\cite{FamaeyBinney2005}, the ``standard'' form
$x/\sqrt{1+x^2}$~\cite{Milgrom1983}, and transcendental forms
such as $1 - e^{-\sqrt{x}}$ all produce acceptable rotation
curves. Famaey and Binney~\cite{FamaeyBinney2005} showed that
the simple form provides the best fit to the Milky Way terminal
velocity curve, but selection by fitting is not derivation. A
derivation of $\mu$ from first principles has remained open for
four decades.

The present letter derives it. The derivation rests on the
coordinate-space geometry of the tetrad frame established in
the companion letter~\cite{Nelson2026a}, which proved that the
tetrad equation $g_{\mu\nu} = \eta_{ab}\,e^a{}_\mu\,e^b{}_\nu$
produces two parity-related bounded tetrahedral regions whose
union is a stellated octahedron. The exchange symmetry between
those two regions provides one of three routes to the
constitutive law developed here.

% ==================================================================

\section{The law}

The gravitational constitutive law is
\begin{equation}
  \boxed{\mu(x) = \frac{x}{x + 1}.}
  \label{eq:law}
\end{equation}
It is the unique function consistent with the exchange symmetry
of the two parity-related sectors of the tetrad equation, under
rational bijectivity on the projective line, as proved in
Section~\ref{sec:route3} and the accompanying
Supplementary Note~3.

Three distinct derivational routes, beginning from different
starting points and using different operative mathematics,
produce this function. All three ultimately identify the same
ratio law, $\mu/(1-\mu) = x$, from which $\mu = x/(x+1)$
follows by normalisation. Throughout, TetA and TetB denote
two local tetrad sectors in the sense of an
interacting-vielbein decomposition~\cite{HR2012a}; the three
routes can be read respectively as a vielbein-structural
derivation, an AQUAL response derivation, and an
exchange-symmetry derivation of the same constitutive law.
The following three sections present them in turn.

% ==================================================================

\section{From the metric}
\label{sec:route1}

The internal metric of the tetrad formalism,
$\eta_{ab} = \mathrm{diag}(-1,+1,+1,+1)$, contains two numbers
readable by inspection: the frame dimension $K = 4$ (the matrix
is $4 \times 4$) and the spatial dimension $D = 3$ (three
positive eigenvalues). Because the tetrad equation is degree~2
in $e$, a third number follows: the root count $n = 2$, whose
geometric realisation as two parity-related bounded regions is
established in the companion letter~\cite{Nelson2026a}.

These three numbers satisfy a closure relation. Two simplices
with $D + 1$ faces each contribute $2(D+1)$ faces; $D$ binary
spatial axes contribute $2^D$ sign combinations. At $D = 3$
the two counts are equal:
\begin{equation}
  2(D + 1) = 2^D = 8.
  \label{eq:closure}
\end{equation}
This equality holds at $D = 3$ as the minimum threshold and
admits no other positive integer solution, since $2^D$ grows
exponentially while $2(D+1)$ grows linearly. The same doubled
count $2K = 8$ appears independently in ghost-free bimetric
gravity~\cite{HR2012a} as the total frame count of two
dynamical tetrad fields, and $K^2 = 16$ as the dimension of
their interaction matrix. The companion
letter~\cite{Nelson2026a} establishes that the four spatial
legs of the tetrad are the four face normals of a regular
tetrahedron under $SO(3)$ isotropy, so frame legs and
tetrahedral faces are the same four directions.

From the same two integers, the tetrad equation's graded
structure can be read. A scalar closure condition
$\det(e) \neq 0$ carries rank~0 on the frame space
(dimension $K^0 = 1$). The frame $e^a{}_\mu$ carries rank~1
(dimension $K^1 = 4$). The bilinear coupling
$\eta_{ab}\,e^a\,e^b$ carries rank~2 (dimension $K^2 = 16$).
Taking the graded dimension of a rank-$r$ ingredient to be
$K^r$, the natural three-term encoding is the coupling
polynomial:
\begin{equation}
  f(K) = 1 + 4 + 16 = 21.
  \label{eq:poly}
\end{equation}
The three ranks contribute independently: the field closure
condition carries no frame index, the frame carries one, and
the bilinear coupling carries two, so their dimensions sum
rather than multiply.

Under spatial parity, a rank-$r$ tensor transforms with sign
$(-1)^r$~\cite{MTW}, so ranks~0 and~2 are even while rank~1
is odd. Because the stellated octahedron has exactly three
structural components (TetA, the octahedral core, and TetB),
these define a natural three-state parity coordinate
$s \in \{+1,\,0,\,-1\}$, with TetA and TetB as the
parity-opposite extrema and the octahedral core as the
symmetric midpoint. The parity-odd contribution enters only
the rank-1 term, and interpolation on three points with odd
symmetry fixes this contribution uniquely as $Ks$:
\begin{equation}
  f(K,s) = 1 + sK + K^2.
  \label{eq:poly_s}
\end{equation}
The polynomial evaluates to $f = 21$ at TetA ($s = +1$),
$f = 17$ at the octahedral core ($s = 0$), and $f = 13$ at
TetB ($s = -1$), with range $\Delta f = 2K = 8$.

To make the discrete state parameter continuous in the
acceleration ratio $x = g/a_0$, a map $s(x)$ is needed. The
map must be rational, to preserve the polynomial's algebraic
structure, and bijective, so that each acceleration corresponds
to exactly one geometric state. A rational bijection on the
projective line is necessarily fractional
linear~\cite{Hartshorne1977}. Three boundary conditions ($s
\to +1$ as $x \to \infty$; $s = 0$ at $x = 1$; $s \to -1$ as
$x \to 0$) determine the unique map $s(x) = (x-1)/(x+1)$.
Composing with the polynomial and normalising by the range
$\Delta f = 8$ yields $\mu(x) = x/(x+1)$.

The full derivation, including the Diophantine verification
table and the consistency check that $K$ and $D$ recover
the principal dimensional and combinatorial properties of
$g_{\mu\nu}$, appears in Supplementary Note~1. This route uses no field
equation and no equilibrium postulates; it reads the
constitutive law from the algebraic content of the metric
and the tetrad equation.

% ==================================================================

\section{From equilibrium}
\label{sec:route2}

The second derivation begins from the AQUAL field equation
itself, which provides the Lagrangian context for modified
dynamics. Because the local tetrad field has four active
coupling ports (the $K = 4$ face directions), the field
response is finitely partitioned across four channels; this
finite partition is what makes a non-trivial TetA/TetB
equilibrium possible.
That capacity partitions between the two coupled sectors at the single interaction
vertex where TetA and TetB meet. The partition reaches
equilibrium when the throughput toward TetA balances the
throughput toward TetB.

At each point, let $\epsilon$ denote the TetA fraction of the
total capacity, so that $1 - \epsilon$ is the TetB fraction.
Two opposing throughputs govern the partition. The
TetA-directed throughput scales as $x^n(1 - \epsilon)$:
stronger acceleration and more available TetB capacity both
push the partition toward TetA. The TetB-directed throughput
scales as $\epsilon$: the larger the TetA fraction, the
stronger the restoring tendency. At equilibrium, these balance
to give
$\epsilon = Kx^n / (1 + Kx^n)$, with $K = \kappa_+/\kappa_-$
the rate ratio. In the small-$x$ limit, $\mu \approx Kx^n$,
so the deep-field scaling $\mu/x \to \text{const}$ requires
$n = 1$. This is the same exponent that produces flat rotation
curves and the baryonic Tully-Fisher relation $v^4 = GMa_0$.

With $n = 1$ established, the equilibrium TetA fraction becomes
$\epsilon = Kx/(1 + Kx)$. The rate ratio $K$ appears only in
the combination $Kx$ and can be absorbed by redefining
$a_0 \to a_0/K$; this is a reparametrisation of the
acceleration scale, not a new degree of freedom. The derivation
therefore fixes the functional form of $\mu$ while leaving the
absolute value of the crossover scale to observation. In
rescaled units,
\begin{equation}
  \mu(x) = \frac{x}{x + 1}.
  \label{eq:mu_eq}
\end{equation}
The equilibrium is globally attracting, with strictly negative
slope in the net rate as a function of $\epsilon$.

The full derivation, including stability analysis and the
corroborating rotation-curve check, appears in
Supplementary Note~2. This route does not
require the coupling polynomial, rank counting, or the
exchange-symmetry derivation; it reads the constitutive law
from the field's equilibrium partition between the two sectors.

% ==================================================================

\section{From exchange symmetry}
\label{sec:route3}

The companion letter~\cite{Nelson2026a} proved that the spatial
sector of the tetrad equation produces exactly two bounded
tetrahedral regions, TetA and TetB, related by spatial parity.
The two sectors couple at the interaction vertex (the
octahedral core where TetA and TetB meet), and the
gravitational field response at each point partitions between
them. The share of that response carried by each sector is its
field capacity at that point.
Let $\mu(x)$ denote the fraction of total field capacity carried
by TetA at dimensionless acceleration $x = g/a_0$, so that
$1 - \mu(x)$ is the fraction carried by TetB. Spatial parity
exchanges the two sectors, inverting the acceleration ratio and
swapping the fractions:
\begin{equation}
  \mu(1/x) = 1 - \mu(x).
  \label{eq:duality}
\end{equation}
The two fractions sum to unity because $\mu$ measures the
relative partition between the two coupled sectors, not the
absolute field strength at a point.

The natural variable for a two-sector partition is the odds
ratio $\Omega = \mu/(1-\mu)$, which measures how many times
more capacity TetA carries than TetB. Under the exchange
symmetry, this ratio inverts when the acceleration inverts:
\begin{equation}
  \Omega(1/x) = \frac{1}{\Omega(x)}.
  \label{eq:odds_duality}
\end{equation}
The endpoint conditions are $\Omega(0) = 0$ (the TetB limit as
$x \to 0$) and $\Omega \to \infty$ as $x \to \infty$
(the TetA limit at large acceleration).

The map $\Omega(x)$ must be rational, to preserve the algebraic
structure of the tetrad equation, and bijective, so that a
single acceleration decodes to exactly one capacity ratio. A
rational bijection on the projective line is necessarily
fractional linear~\cite{Hartshorne1977}, giving the general
form $\Omega(x) = (ax + b)/(cx + d)$. The endpoint condition
$\Omega(0) = 0$ requires $b = 0$, reducing the map to
$\Omega(x) = ax/(cx + d)$. Substituting into the exchange
condition $\Omega(1/x) = 1/\Omega(x)$ and cross-multiplying
gives $a^2 x = (c + dx)(cx + d)$. Comparing coefficients of
$x^2$, $x$, and $x^0$ yields $cd = 0$ and $a^2 = c^2 + d^2$.
If $d = 0$, then $\Omega = a/c$ is a constant, which is not
bijective. If $c = 0$, then $a^2 = d^2$, so $a = \pm d$.
The branch $a = -d$ gives $\Omega(x) = -x$, which is
incompatible with $\Omega > 0$ for $x > 0$; therefore $a = d$
and
\begin{equation}
  \Omega(x) = x.
  \label{eq:odds_result}
\end{equation}
The odds ratio between TetA and TetB is the acceleration ratio
itself. From this, the constitutive law follows by
normalisation:
\begin{equation}
  \mu(x) = \frac{\Omega}{1 + \Omega} = \frac{x}{x + 1}.
  \label{eq:mu_exchange}
\end{equation}

Two properties often treated as independent conditions are in
fact derived consequences. The midpoint value $\mu(1) = 1/2$
follows directly: at the acceleration scale $a_0$, the two
sectors contribute equally, as the exchange symmetry requires.
The complement symmetry $\mu(x) + \mu(1/x) = 1$ is the
duality condition~\eqref{eq:duality} itself, which was the
starting point of the derivation rather than an additional
assumption.

The full derivation, including the explicit coefficient
comparison and the screening identity, appears in
Supplementary Note~3.

% ==================================================================

\section{Why they agree}

Three derivations using different mathematics arrive at the
same function. That convergence has a simple explanation: each
route, through its own formalism, discovers the same algebraic
fact.

In Route~1, the polynomial normalisation $(f - 13)/(21 - f)$
evaluates to $x$ when $f = (21x + 13)/(x+1)$; here $(f - 13)$
measures coupling depth above the TetB floor while $(21 - f)$
measures depth below the TetA ceiling, so their ratio is the
same inter-sector odds that Routes~2 and~3 identify directly.
In Route~2,
the equilibrium condition $\epsilon/(1-\epsilon) = x$ is the
definition of the odds ratio. In Route~3, the exchange symmetry
directly produces $\Omega(x) = x$. All three identify the
capacity ratio between TetA and TetB with the dimensionless
acceleration. The constitutive law is the odds-to-fraction
normalisation of that ratio (the logistic map
$\mu = \Omega/(1+\Omega)$):
\begin{equation}
  \frac{\mu}{1 - \mu} = x
  \qquad \Longrightarrow \qquad
  \mu = \frac{x}{x + 1}.
  \label{eq:convergence}
\end{equation}

The routes are methodologically distinct: Route~1 uses the coupling
polynomial, Route~2 uses the AQUAL field equation, and Route~3
uses the odds ratio. No two share the same primary derivational
mechanism. Their convergence
reflects the fact that a single physical quantity, the
TetA-to-TetB capacity ratio, is accessible through the
algebraic structure of the metric, through the field's
equilibrium kinetics, and through the geometry's exchange
symmetry.

% ==================================================================

\section{Consequences}
\label{sec:consequences}

The constitutive law carries several consequences that have
previously been observed empirically but can now be traced to
the TetA/TetB capacity structure derived in the three routes.

\paragraph{Complement symmetry.}
The relation $\mu(x) + \mu(1/x) = 1$ states that whatever
fraction of field capacity TetA carries at acceleration $x$,
TetB carries at $1/x$. This property was noted by
Milgrom~\cite{Milgrom1983} as a structurally significant
feature of the simple interpolating function. In the present
derivation it is not a property the function happens to
possess; it is the exchange symmetry of the two tetrahedral
sectors, which was the starting point of Route~3.

\paragraph{Baryonic Tully-Fisher relation.}
In the deep-field regime ($x \ll 1$), the constitutive law
reduces to $\mu \approx x$, which corresponds to the
TetB-dominated sector. In spherical symmetry, this limit gives
$v^4 = GMa_0$: the baryonic Tully-Fisher
relation~\cite{McGaugh2000}. One of the strongest empirical
regularities in extragalactic astronomy follows from the
deep-field limit of the constitutive law, which is itself
the TetB sector of the tetrad geometry.

% ==================================================================

\section{Discussion}

The present letter derives the functional form of $\mu$. It
does not derive the value of $a_0$. The acceleration scale
enters as an empirical input that defines the boundary between
the TetA-dominated and TetB-dominated regimes; determining it
from the framework's structural constants is the subject of
ongoing work.

The integers $\{21,\,17,\,13\}$ that appear in Route~1 as the
three evaluation states of the coupling polynomial are developed
in the companion geometry letter~\cite{Nelson2026a}. At the
octahedral core ($s = 0$), the parity-odd rank-1 term
contributes zero, leaving 17 as the coupling depth carried by
the parity-even structure alone. The
constitutive law itself requires none of them: Route~3
derives it from exchange symmetry alone, and Route~2 derives
it from equilibrium kinetics with no reference to the tetrad
geometry. Route~1 contributes what the other two cannot: it
connects the constitutive law to the metric's dimensional
structure and explains why $K = 4$ and $D = 3$ appear in the
polynomial. Each route carries a distinct explanatory role.

The constitutive law $\mu(x) = x/(x+1)$ applies wherever the
AQUAL field equation applies: to any gravitating system in
which the acceleration passes through the scale $a_0$. Its
primary observational arena is galaxy rotation curves, where
it was first identified empirically by
Milgrom~\cite{Milgrom1983} and confirmed by Famaey and
Binney~\cite{FamaeyBinney2005} through fitting to the Milky
Way terminal velocity curve.

The function that has been empirical for forty years now has
a derivation. It is the odds-to-fraction normalisation of
the TetA-to-TetB capacity ratio, recovered through exchange
symmetry, equilibrium closure, and the graded structure of the
metric.

% ==================================================================

\begin{thebibliography}{99}

\bibitem{Milgrom1983}
M.~Milgrom,
\textit{A modification of the Newtonian dynamics as a possible
alternative to the hidden mass hypothesis},
Astrophys.\ J.\ \textbf{270} (1983) 365.

\bibitem{BekensteinMilgrom1984}
J.~D.~Bekenstein and M.~Milgrom,
\textit{Does the missing mass problem signal the breakdown of
Newtonian gravity?}
Astrophys.\ J.\ \textbf{286} (1984) 7.

\bibitem{FamaeyBinney2005}
B.~Famaey and J.~Binney,
\textit{Modified Newtonian dynamics in the Milky Way},
Mon.\ Not.\ R.\ Astron.\ Soc.\ \textbf{363} (2005) 603.

\bibitem{Nelson2026a}
S.~Nelson,
\textit{The Coordinate-Space Geometry of the Tetrad Frame},
Class.\ Quantum Grav.\ (2026), submitted (CQG-115428).

\bibitem{MTW}
C.~W.~Misner, K.~S.~Thorne and J.~A.~Wheeler,
\textit{Gravitation}
(W.~H.~Freeman, San Francisco, 1973).

\bibitem{Hartshorne1977}
R.~Hartshorne,
\textit{Algebraic Geometry}
(Springer, New York, 1977).

\bibitem{HR2012a}
S.~F.~Hassan and R.~A.~Rosen,
\textit{Bimetric gravity from ghost-free massive gravity},
JHEP \textbf{02} (2012) 126.

\bibitem{McGaugh2000}
S.~S.~McGaugh, J.~M.~Schombert, G.~D.~Bothun and
W.~J.~G.~de~Blok,
\textit{The baryonic Tully-Fisher relation},
Astrophys.\ J.\ \textbf{533} (2000) L99.

\end{thebibliography}

\end{document}
